% ============================================================
% CONCLUSIONI
% ============================================================
\section{Conclusioni}

Il progetto ha applicato un insieme di tecniche di verifica software,
manuali e automatiche, a due classi di un fork di Apache Avro:
\texttt{Resolver}, che implementa la logica di risoluzione tra schema
writer e reader, e \texttt{Utf8}, che rappresenta una sequenza di
testo codificata direttamente in byte UTF-8. Per ciascuna classe sono
stati sviluppati test manuali basati su Category-Partition, test
generati automaticamente tramite Randoop e LLM, e la qualità dei test
è stata validata tramite mutation testing (PIT) e una stima di
reliability; infine, entrambe le classi sono state sottoposte a un
refactoring in quattro varianti con contesto crescente.

Un risultato inatteso è emerso già nella fase manuale: durante la
progettazione dei test con Category-Partition per \texttt{Utf8}, è
stato identificato un comportamento potenzialmente anomalo del metodo
\texttt{setByteLength} nei casi di espansione oltre la capacità
attuale del buffer. La verifica diretta del codice sorgente ha
confermato che il metodo non preserva il contenuto dei byte quando il
buffer viene riallocato -- un comportamento non documentato e
controintuitivo, che i soli test automatici non avrebbero portato
alla luce senza un'analisi manuale preliminare.

Sul fronte del random testing, il confronto tra le due classi ha
evidenziato una dipendenza netta dalla complessità del dominio di
input: Randoop ha raggiunto l'81\% di statement coverage su
\texttt{Utf8}, mentre su \texttt{Resolver} è rimasto sostanzialmente
a 0\%, incapace di costruire autonomamente oggetti \texttt{Schema}
semanticamente validi. Questo risultato mostra chiaramente i limiti
del random testing su classi con dipendenze di dominio complesse.

La generazione tramite LLM ha prodotto un risultato metodologicamente
interessante: le suite Zero-Shot, generate con un prompt meno
strutturato, hanno prodotto in media più casi di test e coverage più
alta rispetto alle suite ToT, che utilizzano un prompting più guidato
e dettagliato. Questo suggerisce che un prompt più libero lascia al
modello maggiore spazio per esplorare scenari diversi, mentre la
struttura imposta da ToT tende a far convergere il modello su un
insieme più ristretto e prevedibile di casi. Il risultato è
empiricamente osservabile nei dati raccolti, ma trattandosi di un
singolo esperimento non è generalizzabile: potrebbe essere
influenzato dalla specifica natura delle classi analizzate o dalla
non-determinismo del modello.
